% SEGMENT OLP-0031-S01
% Part:sets-functions-relations
% Chapter: size-of-sets
% Section: pairing

\documentclass[../../../include/open-logic-section]{subfiles}

\begin{document}

\olfileid{sfr}{siz}{pai}
\olsection{ஜோடியாக்கச் சார்புகளும் குறியெண்களும்}

\begin{explain}
காண்டோரின் வளைவழி முறை $\Nat^n$ இன் எண்ணிடத்தக்க தன்மையைப்
பார்வைக்குத் தெளிவாக்குகிறது.
ஆனால் $\Nat^2$ ஐக் காட்டும் நமது அட்டவணையில் கவனம் செலுத்துவோம்.
அட்டவணையிலுள்ள வளைவழியைப் பின்தொடர்ந்து இடங்களை எண்ணினால்,
$\tuple{1,2}$ என்பது $7$ என்ற எண்ணுடன் இணைக்கப்பட்டிருப்பதைக் காணலாம்.
ஆயினும் இதனை மேலும் நேரடியாகக் கணக்கிட முடிந்தால் நன்றாக இருக்கும்.
அதாவது வளைவழிப் பட்டியலாக்கத்தின் \emph{நேர்மாறான}
$g\colon \Nat^2 \to \Nat$ பின்வருமாறு நம்மிடம் இருந்தால் உதவியாக இருக்கும்:
\[
g(\tuple{0,0}) = 0, \;
g(\tuple{0,1}) = 1, \;
g(\tuple{1,0}) = 2, \; \dots,
g(\tuple{1,2}) = 7, \; \dots
\]
நமது பட்டியலாக்கத்தில் $\tuple{n, m}$ சரியாக எந்த இடத்தில்
இடம்பெறும் என்பதைக் கணக்கிட இது உதவும்.

உண்மையில், இரண்டு கவனிப்புகளின் மூலம் $g$ ஐ நேரடியாக வரையறுக்கலாம்.
முதலில், $n$ ஆவது கிடைவரிசையிலும் $m$ ஆவது நெடுவரிசையிலும்
$v$ என்ற மதிப்பு இருந்தால், $(n+1)$ ஆவது கிடைவரிசையிலும்
$(m-1)$ ஆவது நெடுவரிசையிலும் $v + 1$ என்ற மதிப்பு இருக்கும்.
இரண்டாவதாக, நமது பட்டியலாக்கத்தின் முதல் கிடைவரிசை,
$0$, $1$, $3$, $6$ என்று தொடங்கும் முக்கோண எண்களைக் கொண்டுள்ளது.
$k$ ஆவது முக்கோண எண், $< k$ என அமையும் இயல் எண்களின் கூட்டுத்தொகை;
அதை $k(k+1)/2$ எனக் கணக்கிடலாம்.
இந்த இரண்டு கவனிப்புகளையும் இணைத்து, பின்வரும் சார்பைக் கருதுக:
\[
  g(n,m) = \frac{(n+m+1)(n+m)}{2} + n
\]
$g(\tuple{n, m})$ என்பதற்குப் பதிலாகப் பெரும்பாலும்
$g(n, m)$ என்றே எழுதுகிறோம்; அது பார்ப்பதற்கு எளிதாக உள்ளது.
முதலில் $(n+m)^\text{th}$ முக்கோண எண்ணைக் கண்டறிந்து,
பின்னர் அதனுடன் $n$ ஐக் கூட்ட வேண்டும் என்று இது கூறுகிறது.
இது அட்டவணையை நாம் விரும்பியபடியே நிரப்புகிறது.
குறிப்பாக, $\tuple{1, 2}$ என்ற ஜோடி
$\frac{4 \times 3}{2} +
1 = 7$ க்கு அனுப்பப்படுகிறது.

இந்தச் சார்பு $g$, ஜோடிகளைக் கொண்ட ஒரு கணத்தின்
பட்டியலாக்கத்திற்கு \emph{நேர்மாறானது}.
இத்தகைய சார்புகள் \emph{ஜோடியாக்கச் சார்புகள்} எனப்படும்.
\end{explain}

% SEGMENT OLP-0031-S02
\begin{defn}[ஜோடியாக்கச் சார்பு]
$f\colon A \times B \to \Nat$ என்ற சார்பு $f$ ஒன்றுக்கொன்றானதாக இருந்தால்,
அது ஒரு கணக்கியல் \emph{ஜோடியாக்கச் சார்பு} ஆகும்.
$f$ ஆனது $A \times B$ ஐ \emph{குறியாக்குகிறது} என்றும்,
$f(x,y)$ என்பது $\tuple{x,y}$ இன் \emph{குறியெண்} என்றும் கூறுகிறோம்.
\end{defn}

% SEGMENT OLP-0031-S03
\begin{explain}
எடுத்துக்காட்டாக, இயல் எண்களின் ஜோடிகளைக் குறியாக்க
ஜோடியாக்கச் சார்புகளைப் பயன்படுத்தலாம்.
வேறு சொற்களில், உறுப்புகளின் ஒவ்வொரு \emph{ஜோடியையும்}
ஒரே ஒரு \emph{எண்ணால்} குறிக்கலாம்.
ஜோடியாக்கச் சார்பின் நேர்மாறைப் பயன்படுத்தி அந்த எண்ணைக்
\emph{குறிவிலக்கலாம்}; அதாவது அது எந்த ஜோடியைக் குறிக்கிறது
என்பதைக் கண்டறியலாம்.
\end{explain}

% SEGMENT OLP-0031-S04
\begin{prob}
எல்லா எதிர்மமற்ற விகிதமுறு எண்களையும் கொண்ட கணத்தின்
ஒரு பட்டியலாக்கத்தைக் கொடுக்கவும்.
\end{prob}

% SEGMENT OLP-0031-S05
\begin{prob}
$\Rat$ என்பது !!{enumerable} எனக் காட்டுக.
எந்த விகிதமுறு எண்ணையும் $z \in \Int$, $m \in \Nat^+$ ஆகியவற்றைக் கொண்டு
$z/m$ என்ற பின்னமாக எழுதலாம் என்பதை நினைவுகொள்க.
\end{prob}

% SEGMENT OLP-0031-S06
\begin{prob}
$\Bin^*$ இன் ஒரு பட்டியலாக்கத்தை வரையறுக்கவும்.
\end{prob}

% SEGMENT OLP-0031-S07
\begin{prob}
ஒவ்வொரு சாத்தியமான வாய்மை அட்டவணையும் ஒரு வாய்மைச் சார்பை
வெளிப்படுத்துகிறது என்பதை உங்கள் அறிமுகத் தருக்கவியல் பாடத்திலிருந்து
நினைவுகொள்க. வேறு சொற்களில், ஏதேனும் ஓர் $k$ க்கு
$\Bin^k \to \Bin$ என அமையும் சார்புகள் அனைத்துமே வாய்மைச் சார்புகள்.
எல்லா வாய்மைச் சார்புகளையும் கொண்ட கணம் எண்ணிடத்தக்கது என நிறுவுக.
\end{prob}

% SEGMENT OLP-0031-S08
\begin{prob}
!!{enumerable} ஆன ஏதேனும் ஒரு முடிவுறாத கணத்தின் எல்லா முடிவுறு
உட்கணங்களையும் கொண்ட கணம் !!{enumerable} எனக் காட்டுக.
\end{prob}

% SEGMENT OLP-0031-S09
\begin{prob}
% SOURCE CORRECTION OLSIZ-002; see translation/size-notes.tex.
$\Nat$ இன் ஓர் உட்கணம் \emph{இணைமுடிவுறு} எனப்படுவதற்குத்
தேவையானதும் போதுமானதுமான நிபந்தனை,
அது $\Nat$ இன் ஒரு முடிவுறு உட்கணத்திற்கு, இயல் எண்களின்
கணத்திற்குள்ளான நிரப்பியாக இருப்பதாகும்.
அதாவது, $A \subseteq \Nat$ என்பது இணைமுடிவுறு எனப்படுவதற்குத்
தேவையானதும் போதுமானதுமான நிபந்தனை,
$\Nat\setminus A$ முடிவுறுவதாகும்.
$I$ என்ற கணத்தின் !!{element}s, $\Nat$ இன் முடிவுறு மற்றும்
இணைமுடிவுறு உட்கணங்கள் அனைத்துமாக இருக்கட்டும்.
$I$ என்பது !!{enumerable} எனக் காட்டுக.
\end{prob}

% SEGMENT OLP-0031-S10
\begin{prob}
!!{enumerable} கணங்களின் !!{enumerable} சேர்ப்பு
!!{enumerable} எனக் காட்டுக.
அதாவது, $A_1$, $A_2$, \dots{} ஆகியவை கணங்களாகவும்,
ஒவ்வொரு $A_i$ உம் !!{enumerable} ஆகவும் இருந்தால்,
அவை அனைத்தின் சேர்ப்பான $\bigcup_{i=1}^\infty
A_i$ என்பதும் !!{enumerable} ஆகும். [குறிப்பு: இது கடினமானது!]
\end{prob}

% SEGMENT OLP-0031-S11
\begin{prob}
$f \colon A \times B \to \Nat$ என்பது ஏதேனும் ஒரு ஜோடியாக்கச்
சார்பாக இருக்கட்டும். $f$ இன் நேர்மாறு,
$A \times B$ இன் ஒரு பட்டியலாக்கம் எனக் காட்டுக.
\end{prob}

% SEGMENT OLP-0031-S12
\begin{prob}
$\Nat^3$ ஐக் குறியாக்கும் ஒரு சார்பைக் குறிப்பிடுக.
\end{prob}

\end{document}
